\documentclass[11pt]{letter}

\usepackage{geometry}
\geometry{margin=1in}
\usepackage{setspace}
\usepackage{hyperref}

\signature{Decory Edwards}
\address{Decory Edwards \\ Johns Hopkins University \\ Department of Economics}

\begin{document}

\begin{letter}{Hiring Committee \\ Walleye Capital \\ Boston, MA}

\opening{Dear Hiring Committee,}

I am writing to apply for the Quantitative Researcher position with the Quantic team at Walleye Capital. I am currently a PhD candidate in Economics at Johns Hopkins University. My training combines mathematical reasoning, computational modeling, and data driven empirical analysis. I am motivated by problems that require precise logic, careful implementation, and constant validation against data. Walleye’s pragmatic, engineering-driven approach to quantitative investing aligns closely with how I approach research and collaboration.

My research agenda is at the intersection of wealth inequality and macroeconomics. In my job market paper, I build and estimate a structural heterogeneous agent model with returns that differ across households. The core contribution is to show how heterogeneity in returns and income risk jointly shape the wealth distribution and consumption behavior. Relative to closely related work, my framework performs better at matching the lower and middle portions of the wealth distribution. This difference matters quantitatively. Models that focus primarily on returns heterogeneity can fit the top of the wealth distribution closely but tend to generate too much wealth among the bottom half of households. In my framework, income uncertainty generates a precautionary saving motive that produces genuinely low wealth households. This leads to a realistic distribution of marginal propensities to consume and aggregate consumption responses that align with empirical estimates.

A key feature of my research is the heavy use of computation to solve and simulate models that cannot be handled analytically. I write code to solve dynamic programming problems, simulate outcomes under stochastic environments, and evaluate model performance across specifications. I work with large datasets, design hypothesis driven analyses, and focus on robustness and interpretability. This work requires disciplined reasoning and careful validation, skills that map directly to alpha research, portfolio optimization, and automated trading systems.

I am particularly interested in Walleye Capital because of the Quantic team’s emphasis on collaboration, engineering discipline, and real world impact. The opportunity to work alongside researchers, engineers, and traders to develop and deploy fully automated strategies is especially appealing. I value environments where ideas are tested rigorously and where execution quality matters as much as innovation.

My economics training has provided a strong foundation in probability, optimization, and statistical inference. I am comfortable translating theory into working code and analyzing large scale datasets using Python and related tools. I value clear communication and teamwork, both of which are essential when research moves from modeling to production systems.

I am also enthusiastic about the prospect of living and working in Boston and contributing in person to a collaborative research environment. I see this role as an opportunity to apply my academic training in a setting where research outcomes directly affect trading performance.

I would welcome the opportunity to contribute my skills and perspective to Walleye Capital. Thank you for your time and consideration.

\closing{Sincerely,}

\end{letter}
\end{document}
